%% kapa: exam=1 round=31 year=2020 period=1 form=A subject=civil_law subject_ko=민법 seq=34 number=34
%% source: Q-Net 감정평가사 기출문제 (cst00309) artlSeq=5207219 file=extracted/1008/5207219/A형/A형/2020년도 제31회 감정평가사 1차 시험_1교시_A형.pdf
%% conversion: kordoc 4.24.51 -> tools/gen1.py (auto); review=proofread-sonnet; flags=none
\begin{question}{34}
법정지상권에 관한 설명으로 옳지 않은 것을 모두 고른 것은? (다툼이 있으면 판례에 따름)
\begin{framedtext}
ㄱ. X토지에 Y건물의 소유를 위한 법정지상권을 가진 甲의 Y건물을 경매에서 매수한 乙은, 건물철거의 매각조건 등 특별한 사정이 없으면 당연히 법정지상권을 취득한다. \par
ㄴ. X토지를 소유하는 甲이 乙과 함께 그 지상에 Y건물을 신축ㆍ공유하던 중 X토지에 저당권을 설정하였고 그의 실행에 의한 경매에서 丙이 X토지의 소유권을 취득한 경우, Y건물을 위한 법정지상권이 성립하지 않는다. \par
ㄷ. 甲 소유의 X토지와 그 지상건물에 공동저당권이 설정된 후 지상건물을 철거하고 Y건물을 신축하였고 저당권의 실행으로 X토지와 Y건물이 다른 소유자에게 매각된 경우, 특별한 사정이 없으면 Y건물을 위한 법정지상권이 성립한다. \par
ㄹ. X토지에 저당권을 설정한 甲이 저당권자 乙의 동의를 얻어 Y건물을 신축하였으나 저당권의 실행에 의한 경매에서 丙이 X토지의 소유권을 취득한 경우, Y건물을 위한 법정지상권이 성립한다. \par
\end{framedtext}
\begin{choices}
  \item ㄱ, ㄷ
  \item ㄱ, ㄹ
  \item ㄱ, ㄴ, ㄹ
  \item ㄴ, ㄷ, ㄹ
  \item ㄱ, ㄴ, ㄷ, ㄹ
\end{choices}
\end{question}
\answer{4}
